\section{Formal Sketch}
\label{sec:formal}

The verifier of \S\ref{sec:predicate} has a freestanding structural theorem at the primitive layer. Strip the runtime to the three load-bearing structures: a trust store $\mathcal{S}$ with separate issuer-key and kernel-key allowlists, a bilateral envelope $E$ carrying a receipt id, a scope predicate, and two signature slices (one issuer-side, one kernel-side), and a denotational interpreter $\llbracket\cdot\rrbracket$ for the scope predicate against the receipt id. The accept relation is the conjunction
\begin{align*}
  \mathit{accept}(\mathcal{S},E) \,=\;
    & [E.\mathit{issuerSig}.\mathit{kid} \in \mathcal{S}.\mathit{issuerKeys}] \\
    {}\wedge{} & [E.\mathit{kernelSig}.\mathit{kid} \in \mathcal{S}.\mathit{kernelKeys}] \\
    {}\wedge{} & \llbracket E.\mathit{scope}\rrbracket(E.\mathit{receiptId}).
\end{align*}
Signature byte validity is abstracted into trust-store membership: the trust store names keys whose signatures the deployed verifier has separately checked, and the freestanding theorem inspects only the membership and scope conjuncts.

\begin{theorem}[Freestanding accept set]
For every trust store $\mathcal{S}$ and bilateral envelope $E$,
\begin{align*}
  \mathit{accept}(\mathcal{S},E) \;\Longleftrightarrow\;
    & E.\mathit{issuerSig}.\mathit{kid} \in \mathcal{S}.\mathit{issuerKeys} \\
    {}\wedge{} & E.\mathit{kernelSig}.\mathit{kid} \in \mathcal{S}.\mathit{kernelKeys} \\
    {}\wedge{} & \llbracket E.\mathit{scope}\rrbracket(E.\mathit{receiptId}) = \mathit{true}.
\end{align*}
\end{theorem}

The bi-conditional is the load-bearing claim. The forward direction says the verifier returns true only when all three conjuncts hold; the reverse direction says no hidden side channel can sneak admission past any of the three. The proof is one Lean tactic line that unfolds the definition of \emph{accept} and rewrites the resulting Boolean conjunction. The Lean source is self-contained at the module level: it imports the predicate-language type and its denotational interpreter, but takes no dependency on the wider polity model, anchor-quorum machinery, or treaty-intersection construction. A reviewer can audit the file end-to-end at \codepath{formal/lean4/Chio/Chio/Treaty/BilateralAccept.lean}.

Three corollaries follow definitionally. The first, \thm{accept_monotone_in_issuer_store} in \codepath{BilateralAccept.lean}, says accept-set monotonicity in the trust store: if every issuer key in $\mathcal{S}$ is also in $\mathcal{S}'$ and the kernel-key allowlists agree, then $\mathit{accept}(\mathcal{S},E) = \mathit{true}$ implies $\mathit{accept}(\mathcal{S}',E) = \mathit{true}$. Adding a co-signing counterparty's key to the issuer allowlist never causes a previously-admitted envelope to be rejected; the accept set grows monotonically with the trust store. The second, \thm{accept_conj_scope_decompose}, says accept-set conjunction over scope predicates: an envelope whose scope is $p \wedge q$ and which is accepted under $\mathcal{S}$ is also accepted under $\mathcal{S}$ when its scope is rebound to $p$ alone, and also when its scope is rebound to $q$ alone. This is the structural decomposition that lets a verifier evaluate joint admission against a treaty as the intersection of independently named predicates: the accept set under a conjunctive scope equals the intersection of the accept sets under each conjunct. The third, \thm{accept_requires_issuer_key}, is the contrapositive admission test: if the envelope's issuer key id is not in $\mathcal{S}.\mathit{issuerKeys}$, then $\mathit{accept}(\mathcal{S},E) = \mathit{false}$. Rejection implies a trust-store miss; admission requires the issuer key to be in the store as a precondition that is logically prior to the scope check. All three corollaries are proved by unfolding \emph{accept} and rewriting on the freestanding theorem.

\paragraph{Kernel-only axioms.} The module's axiom footprint is minimal. Running Lean's \codepath{#print axioms} against \thm{freestanding_accept_set_theorem} and each of the three corollaries reports only the kernel-level axioms \emph{propext}, \emph{Classical.choice}, and \emph{Quot.sound}: the same trusted base every Mathlib development presupposes. No \emph{sorry} placeholders appear, and the module introduces no custom axioms of its own. The structural claim therefore rests on the same logical foundation as ordinary Lean 4 mathematics, and a reviewer auditing the file only needs to trust the Lean kernel and these three standard axioms; the broader polity machinery, which would introduce predicate-language and constitutional-delta axioms, is left to the companion paper.

The theorem deliberately stops short of three further obligations. It does not prove single-amendment backward refinement: the freestanding statement is over a fixed predicate language, not over amendments that propose a new predicate set, and the trajectory invariant that rules out constitutional ratchets is not in scope at the primitive layer. It does not prove polity-level admission: it speaks only of the verifier's accept set on a bilateral envelope, not of a polity's constitutional admission of a receipt against its own predicate set and citizenship roster. And it does not address the Hart-condition-(a) framing for state-like authority: the construction is a cryptographic primitive over canonical bytes, and the political-theory framing belongs to a different argument.

Those three obligations are discharged separately in the companion polity paper: \thm{amendment_admissible_iff_backward_refinement} carries amendment refinement as a type-level invariant on the enactment constructor; \thm{treaty_admission_iff_predicate_intersection} carries treaty admission as the Boolean intersection of treaty scope, treaty constitution, and both polities' \emph{polityAdmits} relations; and the trajectory-invariant theorem \thm{essential_preserved_chain} rules out a chain of individually valid refinements that collapse to an adversary-only predicate over admitted history. The bilateral-DSSE primitive sits below all three: it is the byte-level admission gate the polity model depends on, with its own structural correctness theorem that does not borrow polity machinery.
